\chapter{System Architecture and Pipeline Design}
\label{ch:arch}

Figure~\ref{fig:aws_diagram} presents the AWS service-level architecture diagram, illustrating the deployed topology including EventBridge scheduling, Lambda compute, SQS decoupling, S3 storage tiers, RDS PostgreSQL, API Gateway, and CloudWatch monitoring.

\begin{figure}[H]
\centering
\includegraphics[width=\textwidth]{aws_techpulse.png}
\caption{AWS architecture diagram.}
\label{fig:aws_diagram}
\end{figure}

\section{Architectural Overview}
\label{sec:arch_overview}

TechPulse has been implemented as a six-layer, cloud-native system deployed on the AWS Free Tier. Each layer is logically independent with well-defined interfaces, enabling modular development, testing, and future extension. The architecture was adapted from the original proposal to operate within Free Tier cost and service constraints, replacing Aurora Serverless with standard RDS PostgreSQL and favoring free-tier LLM APIs while preserving the core hybrid retrieval and MLOps capabilities.

% ---- TikZ: Six-layer architectural stack ----
\begin{figure}[H]
\centering
\resizebox{0.78\textwidth}{!}{%
\begin{tikzpicture}[node distance=0.0cm]

\def\lw{12.0cm}
\def\lh{0.85cm}

\node[rectangle, draw=awsgray!50, fill=awslightblue, minimum width=\lw,
      minimum height=\lh, font=\small\bfseries, text centered,
      rounded corners=2pt]
  (l6) {Layer 6 --- Monitoring \& Interaction \quad \textcolor{awsgray}{\scriptsize CloudWatch, SNS, API GW, React}};

\node[rectangle, draw=awsgray!50, fill=awslightorange, minimum width=\lw,
      minimum height=\lh, font=\small\bfseries, text centered,
      rounded corners=2pt, below=0pt of l6]
  (l5) {Layer 5 --- Generative Synthesis (LLM Backend) \quad \textcolor{awsgray}{\scriptsize Groq / Bedrock / Ollama (auto-fallback)}};

\node[rectangle, draw=awsgray!50, fill=awslightpurple, minimum width=\lw,
      minimum height=\lh, font=\small\bfseries, text centered,
      rounded corners=2pt, below=0pt of l5]
  (l4) {Layer 4 --- Hybrid Retrieval Pipeline \quad \textcolor{awsgray}{\scriptsize filter $\to$ vector + BM25 $\to$ weighted RRF $\to$ cross-encoder}};

\node[rectangle, draw=awsgray!50, fill=awslightgreen, minimum width=\lw,
      minimum height=\lh, font=\small\bfseries, text centered,
      rounded corners=2pt, below=0pt of l4]
  (l3) {Layer 3 --- Vector Storage \& Indexing \quad \textcolor{awsgray}{\scriptsize RDS PostgreSQL + pgvector (HNSW)}};

\node[rectangle, draw=awsgray!50, fill=awslightorange, minimum width=\lw,
      minimum height=\lh, font=\small\bfseries, text centered,
      rounded corners=2pt, below=0pt of l3]
  (l2) {Layer 2 --- Preprocessing \& Embedding \quad \textcolor{awsgray}{\scriptsize 17-step normalize, chunk, SHA-256}};

\node[rectangle, draw=awsgray!50, fill=awslightblue, minimum width=\lw,
      minimum height=\lh, font=\small\bfseries, text centered,
      rounded corners=2pt, below=0pt of l2]
  (l1) {Layer 1 --- Data Ingestion (5 sources) \quad \textcolor{awsgray}{\scriptsize ArXiv, HN, DEV.to, GitHub, RSS}};

% Inter-layer flow arrows (right-hand side, upward ingestion path - perfectly vertical)
\draw[-{Stealth}, thick, color=awsblue!70]
  (l1.east) -| ++(1.8,0.5) |- node[right, font=\scriptsize, color=awsblue!80] {docs} (l2.east);
\draw[-{Stealth}, thick, color=awsgreen!70]
  (l2.east) -| ++(1.8,0.5) |- node[right, font=\scriptsize, color=awsgreen!80] {chunks} (l3.east);
\draw[-{Stealth}, thick, color=awspurple!70]
  (l3.east) -| ++(1.8,0.5) |- node[right, font=\scriptsize, color=awspurple!80] {ranked} (l4.east);
\draw[-{Stealth}, thick, color=awsorange!70]
  (l4.east) -| ++(1.8,0.5) |- node[right, font=\scriptsize, color=awsorange!80] {context} (l5.east);
\draw[-{Stealth}, thick, color=awsblue!70]
  (l5.east) -| ++(1.8,0.5) |- node[right, font=\scriptsize, color=awsblue!80] {response} (l6.east);

% Query flow arrow (left-hand side, downward query path - perfectly vertical)
\draw[-{Stealth}, dashed, thick, color=propgray!70]
  (l6.west) -| ++(-1.8,-0.5) |- node[left, font=\scriptsize, color=propgray!80] {query} (l4.west);

\end{tikzpicture}%
}
\caption{TechPulse six-layer conceptual architecture.}
\label{fig:six_layers}
\end{figure}

\textbf{Layer Architecture Overview:}
TechPulse implements a six-layer, cloud-native stack with well-defined interfaces enabling modular development and testing. The right-hand arrows illustrate the upward data ingestion pipeline: raw documents from five sources (ArXiv, Hacker News, DEV.to, GitHub, RSS) flow through Layer~1 (Ingestion), are normalized and chunked in Layer~2 (Preprocessing), indexed in Layer~3 (Vector Storage with pgvector HNSW), ranked by the Hybrid Retrieval Pipeline in Layer~4 (which combines BM25 and vector similarity via weighted RRF fusion with cross-encoder reranking), synthesized by the LLM backend in Layer~5 with automatic fallback logic (Groq $\to$ Bedrock $\to$ Ollama), and finally exposed via Layer~6 (Monitoring, API Gateway, and React frontend). The left-hand dashed arrow represents the inbound query path from Layer~6 to Layer~4 for retrieval initiation. This layered design aligns with the AWS Well-Architected Framework principles while maintaining Free Tier cost constraints.

\section{Alignment with the AWS Well-Architected Framework}
\label{sec:well_architected}

The TechPulse architecture is designed in explicit alignment with the AWS Well-Architected Framework (AWS, 2023). Table~\ref{tab:well_arch} summarizes how each pillar is addressed within the implemented system.

\begin{table}[H]
\caption{TechPulse Alignment with the AWS Well-Architected Framework Six Pillars.}
\label{tab:well_arch}
\centering
\begin{tabular}{p{3.2cm}p{9.3cm}}
\toprule
\textbf{Pillar} & \textbf{TechPulse Implementation} \\
\midrule
\textbf{Operational Excellence} & Automated ingestion via EventBridge (6-hour cloud cycle; per-source intervals in Table~\ref{tab:data_sources}); CI/CD pipeline using GitHub Actions and SAM; CloudWatch logging, custom metrics, and health-check Lambda (5-minute interval) \\[4pt]
\textbf{Security} & Controlled API access via Amazon API Gateway with CORS; rate limiting via slowapi; all services managed within AWS default VPC; no direct model invocation from the frontend \\[4pt]
\textbf{Reliability} & Decoupled ingestion and embedding pipelines via SQS; Dead Letter Queue for failure recovery and replay with 14-day retention; automatic LLM backend fallback chain (zero per-backend retries; fallback to next tier provides redundancy within the 30\,s API Gateway timeout); request-level deadline budget with adaptive cross-encoder skip \\[4pt]
\textbf{Performance Efficiency} & Serverless Lambda scaling; lightweight MiniLM embedding model (384-dim); metadata pre-filtering reduces vector search scope; HNSW index for approximate nearest-neighbor search \\[4pt]
\textbf{Cost Optimization} & Serverless-only execution (pay-per-invocation); programmatic budget guard via Cost Explorer; Groq free-tier inference as primary backend (500k tokens/day); Amazon Bedrock available in Free Tier; db.t3.micro RDS instance \\[4pt]
\textbf{Sustainability} & Event-driven processing avoids idle compute; SQS decoupling prevents over-provisioning; right-sized Lambda memory allocations (512MB--2048MB depending on workload); 5-minute warmup pings to keep critical Lambda functions warm \\
\bottomrule
\end{tabular}
\end{table}

\section{Data Sources}
\label{sec:data_sources}

The original proposal specified two data sources (ArXiv and Hacker News). During implementation, the data source coverage was expanded to five heterogeneous sources to improve corpus diversity, topical breadth, and temporal coverage. Table~\ref{tab:data_sources} summarizes the implemented data sources.

\begin{table}[H]
\caption{Implemented Data Sources and Roles.}
\label{tab:data_sources}
\centering
\begin{tabular}{p{3.0cm}p{2.5cm}p{4.5cm}p{2.5cm}}
\toprule
\textbf{Source} & \textbf{Type} & \textbf{Role} & \textbf{Frequency} \\
\midrule
ArXiv API & Batch & Scholarly grounding corpus (cs.AI, cs.LG, cs.CL, cs.IR, cs.SE) & Every 12 hours \\
Hacker News API & Near-real-time & Community-curated tech discussions & Every 30 minutes \\
DEV.to API & Near-real-time & Practitioner tutorials and industry articles (21 tags) & Every 30 minutes \\
GitHub Trending & Batch & Open-source project intelligence (10 topic queries) & Every 30 minutes \\
RSS/Atom Feeds & Batch & Curated technology journalism (10 feeds) & Every 30 minutes \\
\bottomrule
\end{tabular}
\end{table}

The frequencies in Table~\ref{tab:data_sources} reflect the local scheduler (\texttt{scheduler.py}), which provides fine-grained per-source intervals. On AWS, a single EventBridge rule triggers the ingestion Lambda every 6 hours, executing all five sources in one invocation.

The combination provides complementary temporal and authority profiles: ArXiv contributes rigorously authored research with stable identifiers; Hacker News contributes practitioner signals and community-curated trending discussions; DEV.to provides developer-focused tutorials and implementation insights; GitHub Trending captures emerging open-source tools and frameworks; and curated RSS feeds from outlets including Ars Technica, TechCrunch, Wired, IEEE Spectrum, and MIT News provide technology journalism coverage. Incoming source values are validated at the API boundary against \texttt{VALID\_SOURCES = \{arxiv, hn, github, devto, rss\}}. The preprocessing pipeline also applies approximately 110 ecosystem-specific domain stopwords (covering terms such as ``model'', ``framework'', ``training'', ``evaluation'') that are too prevalent in the technology corpus to carry discriminative retrieval or keyword-extraction value.

\section{Data Engineering Pipeline}
\label{sec:pipeline}

The pipeline has been implemented around four design principles: determinism (identical input produces identical output), idempotency (re-running does not create duplicate records), observability (all state transitions are logged), and failure resilience (failures are isolated and recoverable).

% ---- TikZ: Full data engineering pipeline ----
\begin{figure}[H]
\centering
\resizebox{\textwidth}{!}{%
\begin{tikzpicture}[node distance=0.4cm and 0.5cm, font=\small]

% ---- Sources ----
\node[compnode=awsblue] (arxiv) {ArXiv API\\(12 h)};
\node[compnode=awsblue, below=0.25cm of arxiv] (hn) {Hacker News\\(30 min)};
\node[compnode=awsblue, below=0.25cm of hn] (devto) {DEV.to\\(30 min)};
\node[compnode=awsblue, below=0.25cm of devto] (github) {GitHub\\(30 min)};
\node[compnode=awsblue, below=0.25cm of github] (rss) {RSS Feeds\\(30 min)};

% ---- EventBridge ----
\node[compnode=awsorange, right=0.8cm of hn, yshift=-0.55cm] (eb) {EventBridge\\Scheduler};

% ---- Ingestion Lambda ----
\node[compnode=awsorange, right=0.7cm of eb] (inlambda) {Ingestion\\Lambda};

% ---- S3 Raw ----
\node[compnode=awsgreen, above right=0.1cm and 0.7cm of inlambda] (s3raw) {S3 Raw\\(5 source prefixes)};

% ---- SQS ----
\node[compnode=awsgray, right=0.7cm of inlambda, yshift=-0.55cm] (sqs) {SQS\\Queue};

% ---- DLQ ----
\node[compnode=propred!80, below=0.35cm of sqs] (dlq) {Dead Letter\\Queue};

% ---- Preprocess + Embed Lambda ----
\node[compnode=awsorange, right=0.7cm of sqs] (pelambda) {Preprocess \&\\Embed Lambda\\(MiniLM)};

% ---- S3 Processed ----
\node[compnode=awsgreen, above right=0.1cm and 0.7cm of pelambda] (s3proc) {S3 Processed\\(/processed)};

% ---- RDS ----
\node[compnode=awspurple, right=0.7cm of pelambda, yshift=-0.55cm] (rds) {RDS PG\\+ pgvector};

% ---- Retrieval / LLM ----
\node[compnode=awsorange, right=0.7cm of rds] (rag) {RAG\\Orchestrator};
\node[compnode=awsblue, right=0.7cm of rag] (llm) {LLM Backend\\(Groq/Bedrock)};

% ---- API GW + React ----
\node[compnode=awsgreen, right=0.7cm of llm] (apigw) {API\\Gateway};
\node[compnode=awsblue, right=0.5cm of apigw] (react) {React\\Frontend};

% ---- Observability ----
\node[compnode=propgray, below=0.8cm of rag, xshift=0.5cm] (cw) {CloudWatch\\Metrics/Logs};
\node[compnode=propgray, right=0.6cm of cw] (sns) {SNS\\Alerts};
\node[compnode=propgray, left=0.6cm of cw] (hc) {Health Check\\Lambda};

% Arrows: sources -> eb -> inlambda
\draw[parrow] (arxiv.east) -- ++(0.25,0) |- (eb.west);
\draw[parrow] (hn.east)    -- ++(0.25,0) |- (eb.west);
\draw[parrow] (devto.east) -- ++(0.25,0) |- (eb.west);
\draw[parrow] (github.east) -- ++(0.25,0) |- (eb.west);
\draw[parrow] (rss.east)   -- ++(0.25,0) |- (eb.west);
\draw[parrow] (eb) -- (inlambda);
\draw[parrow] (inlambda.east) -- ++(0.25,0) |- (s3raw.west);
\draw[parrow] (inlambda) -- (sqs);
\draw[parrow, dashed, color=propred] (sqs.south) -- (dlq.north)
  node[midway, right, font=\scriptsize, color=propred] {failures};
\draw[parrow] (sqs) -- (pelambda);
\draw[parrow] (pelambda.east) -- ++(0.25,0) |- (s3proc.west);
\draw[parrow] (pelambda) -- (rds);
\draw[parrow] (rds) -- (rag);
\draw[parrow] (rag) -- (llm);
\draw[parrow] (llm) -- (apigw);
\draw[parrow] (apigw) -- (react);

% Observability arrows
\draw[parrow, dashed, color=propgray] (rag.south) -- (cw.north);
\draw[parrow, dashed, color=propgray] (cw.east) -- (sns.west);
\draw[parrow, dashed, color=propgray] (hc.east) -- (cw.west);

% Layer labels
\node[below=0.5cm of devto, font=\scriptsize\bfseries, color=awsblue, xshift=0.2cm]
  {Sources};
\node[below=0.5cm of eb, font=\scriptsize\bfseries, color=awsorange]
  {Trigger};
\node[below=1.2cm of sqs, font=\scriptsize\bfseries, color=awsorange, xshift=0.8cm]
  {Processing};
\node[below=0.5cm of rds, font=\scriptsize\bfseries, color=awspurple]
  {Storage};
\node[below=0.5cm of rag, font=\scriptsize\bfseries, color=awsorange, xshift=0.3cm]
  {Serving};
\node[below=0.3cm of cw, font=\scriptsize\bfseries, color=propgray]
  {Observability};

\end{tikzpicture}%
}
\caption{End-to-end data engineering and serving pipeline}
\label{fig:pipeline}
\end{figure}

\subsection{Ingestion and Decoupling}

Ingestion has been implemented through AWS Lambda functions orchestrated by Amazon EventBridge, supporting both scheduled and event-driven execution. Each of the five ingesters (ArXiv, Hacker News, DEV.to, GitHub, RSS) operates independently with source-specific data extraction logic. Raw artifacts are persisted in Amazon S3 under a structured prefix hierarchy (\texttt{/raw/arxiv/}, \texttt{/raw/hn/}, \texttt{/raw/devto/}, \texttt{/raw/github/}, \texttt{/raw/rss/}, \texttt{/processed/}) to preserve transformation lineage. The ingestion pipeline is decoupled from the embedding pipeline via Amazon SQS, providing backpressure control, failure isolation, and independent horizontal scaling. Messages that fail processing after three retry attempts are routed to a Dead Letter Queue with 14-day retention for controlled replay.

All five ingesters share a common architectural pattern: (1) fetch data from API or feed; (2) extract and normalize fields; (3) compute SHA-256 content hash over canonicalized fields; (4) insert into the \texttt{documents} table with \texttt{ON CONFLICT (content\_hash) DO NOTHING} for idempotent deduplication; (5) write raw JSON to S3; and (6) enqueue document ID to SQS for downstream processing. After indexing, a RAG orchestrator coordinates retrieval and generation by interfacing with the LLM backend fallback chain, as illustrated in Figure~\ref{fig:pipeline} and detailed in Section~\ref{sec:llm_fallback}.

\subsection{Document State Lifecycle}

Each document progresses through a four-state deterministic lifecycle to guarantee idempotency and traceability. Table~\ref{tab:doc_states} defines these states.

\begin{table}[H]
\caption{Document State Transition Model.}
\label{tab:doc_states}
\centering
\begin{tabular}{p{2.0cm}p{3.8cm}p{3.0cm}p{3.0cm}}
\toprule
\textbf{State} & \textbf{Description} & \textbf{Storage} & \textbf{Trigger} \\
\midrule
RAW & Document ingested from API & Amazon S3 & Successful API fetch \\
PROCESSED & Cleaned, normalized, chunked & S3 (\texttt{/processed}) & Preprocessing complete \\
EMBEDDED & Vector representation generated & RDS (pgvector) & Embedding computed \\
INDEXED & Available for retrieval & RDS metadata table & Vector committed \\
\bottomrule
\end{tabular}
\end{table}

Figure~\ref{fig:state_machine} illustrates the document state machine, showing the trigger and storage destination for each transition.

% ---- TikZ: Document State Machine ----
\begin{figure}[H]
\centering
\resizebox{\textwidth}{!}{%
\begin{tikzpicture}[node distance=0.5cm and 2.0cm, font=\small,
  snode/.style={rectangle, draw=awsblue, fill=awslightblue, rounded corners=4pt,
                minimum width=2.2cm, minimum height=0.7cm, align=center, font=\small\bfseries}]

\node[snode] (raw)      {RAW};
\node[snode, right=2.5cm of raw]   (proc)  {PROCESSED};
\node[snode, right=2.5cm of proc]  (emb)   {EMBEDDED};
\node[snode, right=2.5cm of emb]   (idx)   {INDEXED};

\draw[-{Stealth}, thick]
  (raw)  -- node[above, font=\scriptsize] {Preprocess}
             node[below, font=\scriptsize, color=awsgreen!70] {S3 /processed} (proc);
\draw[-{Stealth}, thick]
  (proc) -- node[above, font=\scriptsize] {Embed (MiniLM)}
             node[below, font=\scriptsize, color=awspurple!70] {RDS pgvector} (emb);
\draw[-{Stealth}, thick]
  (emb)  -- node[above, font=\scriptsize] {pgvector commit}
             node[below, font=\scriptsize, color=awsorange!70] {RDS metadata} (idx);

% Entry arrow
\draw[-{Stealth}, thick, color=awsblue!70]
  ([xshift=-1.4cm]raw.west) -- node[above, font=\scriptsize] {Ingestion} (raw.west);

% S3 annotation for RAW
\node[font=\scriptsize, color=awsgreen!70, below=0.6cm of raw] {S3 /raw/\{source\}};

\end{tikzpicture}%
}
\caption{Document state machine with four deterministic states.}
\label{fig:state_machine}
\end{figure}

To guarantee idempotency, canonicalized content fields are hashed using SHA-256 (as described in Section~\ref{sec:innovation}) and deduplicated via PostgreSQL \texttt{ON CONFLICT} clauses, ensuring consistent corpus integrity across experimental runs.

\subsection{Preprocessing and Chunking}

Textual preprocessing has been implemented as a comprehensive 17-step normalization pipeline in the \texttt{normalize\_text()} function. Table~\ref{tab:preprocess_steps} summarizes the grouped operations applied in sequence.

\begin{table}[H]
\caption{Preprocessing Pipeline: 17-Step Text Normalization.}
\label{tab:preprocess_steps}
\centering
\begin{tabular}{p{3.5cm}p{1.0cm}p{7.5cm}}
\toprule
\textbf{Category} & \textbf{Steps} & \textbf{Operations} \\
\midrule
HTML Cleanup & 1--3 & Entity unescaping, comment removal, tag stripping \\[2pt]
Code \& Markdown & 4--6 & Fenced/inline code removal, image/link normalization \\[2pt]
URL Removal & 7 & HTTP/HTTPS/FTP URL stripping \\[2pt]
Formatting Cleanup & 8--12 & Heading, bold/italic, strikethrough, rule, blockquote markers \\[2pt]
Unicode Normalization & 13--14 & Emoji/pictographic removal (6 Unicode ranges), control characters \\[2pt]
Whitespace Normalization & 15--17 & Punctuation collapse, NBSP normalization, whitespace trimming \\
\bottomrule
\end{tabular}
\end{table}

Following normalization, documents are segmented using token-based chunking with the \texttt{tiktoken} tokenizer (\texttt{cl100k\_base} encoding). The implemented configuration uses a 500-token window with 50-token overlap. The 500-token window was selected to fit comfortably within the top-$k$ context budget: with $k = 5$ chunks, the total retrieved context consumes approximately 2,500 tokens---well within the 128k-token window of Groq-hosted Llama 3.1 8B Instant and the 128k-token window of Amazon Nova Micro (Bedrock), while providing sufficient per-chunk semantic coherence. The 50-token (10\%) overlap prevents important information at chunk boundaries from being split across non-adjacent retrieval candidates.

\section{Embedding and Indexing}
\label{sec:embedding}

Semantic representations are generated using the \texttt{all-MiniLM-L6-v2} sentence-transformer model (Wang et al., 2020), producing 384-dimensional embeddings. The model is loaded as a lazy-initialized singleton to minimize memory overhead under serverless execution constraints. Embeddings are stored in RDS PostgreSQL using the pgvector extension with an HNSW (Hierarchical Navigable Small World) index configured with $m = 16$ and $\text{ef\_construction} = 64$, enabling efficient approximate nearest-neighbor search with cosine distance. The parameters $m = 16$ (maximum number of bidirectional links per node) and $\text{ef\_construction} = 64$ (size of the dynamic candidate list during index construction) follow the pgvector recommended defaults for datasets up to $\sim$100k vectors, providing a favorable balance between index build time, memory footprint, and recall (pgvector documentation recommends $m \in [8, 64]$ and $\text{ef\_construction} \geq 2 \times m$).

\section{Hybrid Retrieval Pipeline Architecture}
\label{sec:retrieval_detail}

The implemented hybrid retrieval pipeline combines vector similarity, full-text (BM25) search, and learning-to-rank fusion. Figure~\ref{fig:retrieval_pipeline} illustrates the four-stage pipeline architecture with per-stage timing estimates and implementation details.

% ---- TikZ: Simplified Hybrid Retrieval Pipeline ----
\begin{figure}[H]
\centering
\resizebox{0.55\textwidth}{!}{%
\begin{tikzpicture}[node distance=1.2cm, font=\small, every text node part/.style={align=center}]

% Stage nodes
\node[rectangle, draw=awsblue, fill=awslightblue, rounded corners=4pt, minimum width=2.0cm, minimum height=0.6cm] (stage0) {Query};

\node[rectangle, draw=awsorange, fill=awslightorange, rounded corners=4pt, minimum width=2.2cm, minimum height=0.7cm, below=of stage0] (stage1) {Metadata\\Filter};

\node[rectangle, draw=awspurple, fill=awslightpurple, rounded corners=4pt, minimum width=2.0cm, minimum height=0.6cm, below=of stage1, xshift=-2.2cm] (stage2a) {Vector ANN};

\node[rectangle, draw=awsgreen, fill=awslightgreen, rounded corners=4pt, minimum width=2.0cm, minimum height=0.6cm, below=of stage1, xshift=2.2cm] (stage2b) {BM25};

\node[rectangle, draw=awsorange!70, fill=awsorange!20, rounded corners=4pt, minimum width=2.5cm, minimum height=0.7cm, below=1.0cm of stage1, xshift=0cm] (stage3) {RRF Fusion};

\node[rectangle, draw=propblue!70, fill=propblue!20, rounded corners=4pt, minimum width=2.5cm, minimum height=0.7cm, below=of stage3] (stage4) {Cross-Encoder};

\node[rectangle, draw=awsblue, fill=awslightblue, rounded corners=4pt, minimum width=2.0cm, minimum height=0.6cm, below=of stage4] (output) {Top-5 Results};

% Arrows (very thick for visibility at small scale)
\draw[-{Stealth}, line width=2pt] (stage0) -- (stage1) node[midway, right, font=\scriptsize] {normalize};
\draw[-{Stealth}, line width=2pt] (stage1) -- (stage2a) node[midway, above left, font=\scriptsize] {parallel};
\draw[-{Stealth}, line width=2pt] (stage1) -- (stage2b) node[midway, above right, font=\scriptsize] {parallel};
\draw[-{Stealth}, line width=2pt] (stage2a) -| (stage3.west);
\draw[-{Stealth}, line width=2pt] (stage2b) -| (stage3.east);
\draw[-{Stealth}, line width=2pt] (stage3) -- (stage4);
\draw[-{Stealth}, line width=2pt] (stage4) -- (output);

% Timing labels
\node[font=\scriptsize, below=0.15cm of stage1, xshift=-1.3cm] {50ms};
\node[font=\scriptsize, below=0.15cm of stage2a] {250ms};
\node[font=\scriptsize, below=0.15cm of stage2b] {150ms};
\node[font=\scriptsize, below=0.15cm of stage3] {50ms};
\node[font=\scriptsize, below=0.15cm of stage4] {350ms};

\end{tikzpicture}%
}
\caption{Hybrid retrieval pipeline architecture with four stages}
\label{fig:retrieval_pipeline}
\end{figure}

\subsubsection{Request Deadline and Adaptive Cross-Encoder Skip}

The API Gateway HTTP API imposes a hard 30-second integration timeout. To ensure the full RAG pipeline (retrieval + LLM generation) completes within this budget, a 25-second request deadline is set at the start of each query (\texttt{\_request\_deadline = time.time() + 25}). This deadline is passed to the hybrid retrieval function, which monitors remaining time before invoking the cross-encoder reranking stage. If fewer than 5 seconds remain when cross-encoder execution would begin, the stage is skipped entirely and the RRF-fused scores are returned directly---preserving enough time for LLM generation. This adaptive behavior prevents gateway timeouts on queries with unusually slow retrieval (e.g., large BM25 corpus scans or PostgreSQL connection pool contention).

\subsubsection{Score Normalization}

Two normalization strategies are applied depending on the pipeline path. When the cross-encoder executes, raw logit scores are transformed via \textbf{sigmoid normalization} ($\sigma(x) = 1/(1 + e^{-x})$) to map the unbounded logit range ($-\infty, +\infty$) into $[0, 1]$ for consistent display and ranking. When the cross-encoder is skipped (deadline or failure), \textbf{min-max normalization} is applied to the RRF fusion scores to produce the same $[0, 1]$ output range. Both strategies ensure that source confidence scores displayed to the user are interpretable percentages regardless of the retrieval path taken.

\subsubsection{Connection Pool Resilience}

The retrieval layer wraps database queries in a retry mechanism that catches \texttt{psycopg2.InterfaceError}, \texttt{psycopg2.OperationalError}, and \texttt{psycopg2.pool.PoolError}. On transient connection failures (e.g., pool exhaustion under concurrent requests), the system discards the failed connection, acquires a fresh one from the pool, and retries the query once before propagating the error. This prevents intermittent PostgreSQL connection drops from surfacing as user-visible failures.

\section{Storage Architecture (Medallion Layout)}
\label{sec:storage}

\begin{table}[H]
\caption{Storage Layers and Responsibilities (Medallion Architecture).}
\label{tab:storage_layers}
\centering
\begin{tabular}{p{2.8cm}p{5.8cm}p{4.0cm}}
\toprule
\textbf{Layer} & \textbf{Technology} & \textbf{Purpose} \\
\midrule
Raw (Bronze) & Amazon S3 (\texttt{/raw/\{source\}/}) & Store unprocessed API responses \\
Processed (Silver) & Amazon S3 (\texttt{/processed/}) & Store cleaned and chunked documents \\
Vector (Gold) & RDS PostgreSQL + pgvector & Store embeddings and structured metadata \\
\bottomrule
\end{tabular}
\end{table}

\section{Database Schema}
\label{sec:schema}

The PostgreSQL database schema implements three core tables. The \texttt{documents} table stores ingested records with fields for source, title, content, URL, publication date, SHA-256 content hash (unique constraint), document state, and timestamps. The \texttt{chunks} table stores preprocessed text segments with their 384-dimensional vector embeddings, linked to parent documents via foreign key. The \texttt{drift\_baselines} table persists drift detection history including mean similarity, standard deviation, probe count, and alert status.

Four database indexes support efficient retrieval: an HNSW vector index on chunk embeddings (\texttt{vector\_cosine\_ops}, $m=16$, $\text{ef\_construction}=64$), and B-tree indexes on \texttt{published\_at} (descending), \texttt{source}, and \texttt{state} for metadata filtering performance.

Figure~\ref{fig:er_diagram} presents the entity--relationship diagram of the PostgreSQL schema.

% ---- TikZ: Entity-Relationship Diagram ----
\begin{figure}[H]
\centering
\resizebox{0.82\textwidth}{!}{%
\begin{tikzpicture}[node distance=0.4cm and 3.2cm, font=\scriptsize]

% documents table
\node[rectangle, draw=awsblue, fill=awslightblue, rounded corners=3pt,
      minimum width=4.2cm, text width=4.0cm, align=left, inner sep=6pt]
  (docs) {
  \textbf{\small documents}\\
  \underline{id} \hfill SERIAL (PK)\\
  source \hfill VARCHAR(20)\\
  title \hfill TEXT\\
  content \hfill TEXT\\
  url \hfill TEXT\\
  published\_at \hfill DATE\\
  content\_hash \hfill CHAR(64) (UNIQUE)\\
  state \hfill VARCHAR(10)\\
  created\_at \hfill TIMESTAMP\\
  \textit{Index: B-tree (published\_at)}\\
  \textit{Index: B-tree (source)}\\
  \textit{Index: B-tree (state)}
  };

% chunks table
\node[rectangle, draw=awspurple, fill=awslightpurple, rounded corners=3pt,
      minimum width=4.2cm, text width=4.0cm, align=left, inner sep=6pt,
      right=3.0cm of docs]
  (chunks) {
  \textbf{\small chunks}\\
  \underline{id} \hfill SERIAL (PK)\\
  document\_id \hfill INTEGER (FK)\\
  chunk\_index \hfill INTEGER\\
  chunk\_text \hfill TEXT\\
  embedding \hfill VECTOR(384)\\
  created\_at \hfill TIMESTAMP\\
  \textit{Index: HNSW (embedding)}
  };

% drift_baselines table
\node[rectangle, draw=awsorange, fill=awslightorange, rounded corners=3pt,
      minimum width=4.2cm, text width=4.0cm, align=left, inner sep=6pt,
      below=1.4cm of docs]
  (drift) {
  \textbf{\small drift\_baselines}\\
  \underline{id} \hfill SERIAL (PK)\\
  run\_date \hfill TIMESTAMP\\
  mean\_similarity \hfill FLOAT\\
  std\_similarity \hfill FLOAT\\
  num\_probes \hfill INTEGER\\
  alert\_triggered \hfill BOOLEAN
  };

% Relationship: documents -> chunks (1 to many)
\draw[-{Stealth}, thick, color=awspurple!70]
  (docs.east) -- node[above, font=\scriptsize, color=awspurple!80] {1} (chunks.west)
  node[pos=0.85, above, font=\scriptsize, color=awspurple!80] {N};

% Indexes are now inside the documents node above

\end{tikzpicture}%
}
\caption{Entity--relationship diagram of the TechPulse PostgreSQL schema.}
\label{fig:er_diagram}
\end{figure}

\section{Prompt Construction and Grounded Generation}
\label{sec:prompt_ch4}

Retrieved top-$k$ documents are concatenated into a structured prompt template that enforces the three-layer hallucination verification strategy (Section~\ref{sec:innovation}): (1) a \textbf{System Instruction} mandating evidence-grounded responses with \texttt{[Source N]} citation format; (2) a \textbf{Context Block} with retrieved excerpts, source identifiers, titles, and dates; (3) the \textbf{User Question}; and (4) a \textbf{Citation Instruction}.

\begin{lstlisting}[caption={Implemented Structured Generation Prompt Template.}, label={lst:prompt}]
SYSTEM:
Answer using the provided context documents as your
primary evidence. For each specific claim, cite the
source in [Source N] format. Do NOT speculate beyond
what the documents support.

CONTEXT:
[Source 1] Title: ... | Date: ... | Source: arxiv
Excerpt text...

[Source 2] Title: ... | Date: ... | Source: hn
Excerpt text...

QUESTION:
{User Query}

INSTRUCTION:
Provide a concise answer and cite supporting sources.
\end{lstlisting}

If retrieval returns insufficient evidence or if the citation grounding check fails, the system returns an abstention response rather than generating unsupported content. Three distinct fallback messages are defined: (1)~\texttt{ABSTENTION\_MESSAGE} when grounding verification fails (``The system could not produce a sufficiently grounded response\ldots''); (2)~\texttt{BUDGET\_HALT\_MESSAGE} when the runtime budget guard (\texttt{\_check\_budget\_exceeded()}) detects that monthly Cost Explorer spend has reached the configured threshold, halting LLM invocation while still returning retrieval results; and (3)~\texttt{LLM\_FALLBACK\_MESSAGE} when all backends in the fallback chain are exhausted (``The language model is temporarily unavailable\ldots'').

\section{LLM Backend Fallback Chain}
\label{sec:llm_fallback}

To ensure generation availability across diverse deployment environments, TechPulse implements an automatic LLM backend fallback chain. When the primary backend fails (e.g., due to service unavailability, rate limiting, or network errors), the system automatically cascades to the next available backend without manual intervention.

\begin{table}[H]
\caption{Implemented LLM Backend Fallback Hierarchy.}
\label{tab:llm_backends}
\centering
\begin{tabular}{p{0.5cm}p{3.5cm}p{4.0cm}p{4.0cm}}
\toprule
\textbf{Tier} & \textbf{Backend} & \textbf{Fallback Chain} & \textbf{Notes} \\
\midrule
1 & Groq (Llama 3.1 8B Instant) & $\to$ Bedrock $\to$ Ollama $\to$ HuggingFace & Primary backend; free-tier API; 500k tokens/day \\[4pt]
2 & Amazon Bedrock (Nova Micro) & $\to$ Groq $\to$ Ollama $\to$ HuggingFace & AWS-native fallback; Free Tier eligible \\[4pt]
3 & Ollama (local) & $\to$ Groq $\to$ HuggingFace & OpenAI-compatible API; local inference \\[4pt]
4 & HuggingFace Inference API & $\to$ Groq $\to$ Ollama & Fallback option; Mistral-7B-Instruct \\
\bottomrule
\end{tabular}
\end{table}

All four backends produce responses compatible with the same structured prompt template, ensuring that evaluation results remain comparable across backends. Each backend is configured with zero per-backend retries; instead, the fallback chain itself provides redundancy by cascading immediately to the next tier on failure. This design is driven by the API Gateway HTTP API's hard 30-second integration timeout: a single Groq timeout (15\,s) plus one retry (15\,s) would already exhaust the budget before reaching any fallback. By eliminating retries and relying on the fallback chain, the system maximizes the probability of returning a generated answer within the timeout window. The Bedrock backend in particular enforces its timeout via \texttt{botocore.Config(read\_timeout=15, connect\_timeout=5, retries=\{"max\_attempts": 0\})}, ensuring that a slow AWS endpoint cannot silently consume the remaining deadline. If all backends fail, the system degrades gracefully by returning a retrieval-only summary constructed from the top-$k$ documents without generative expansion.

% ---- TikZ: LLM Backend Fallback Chain Flowchart ----
\begin{figure}[H]
\centering
\resizebox{\textwidth}{!}{%
\begin{tikzpicture}[node distance=0.6cm and 2.0cm, font=\small, every path/.style={thick},
                     tier/.style={rectangle, rounded corners=4pt, minimum width=2.8cm,
                                  minimum height=0.7cm, align=center, text width=2.6cm}]

% Entry node
\node[rectangle, draw=awsblue, fill=awslightblue, rounded corners=4pt,
      minimum width=2.5cm, minimum height=0.7cm]
  (query) {Query Received};

% Tier 1: Groq
\node[tier, draw=awsorange, fill=awslightorange, below=0.8cm of query]
  (groq) {1. Groq (Primary)\\Llama 3.1 8B};

% Retry logic (implicit, shown in label)
\draw[parrow] (query) -- (groq);
\draw[-{Stealth}, thick, dashed, color=awsorange!70]
  ([xshift=1.4cm]groq.east) arc (0:-90:0.8cm)
  node[right, font=\scriptsize, color=awsorange!70] {fallback};

% Tier 2: Bedrock (fork on Groq failure)
\node[tier, draw=awspurple, fill=awslightpurple, right=2.0cm of groq]
  (bedrock) {2. Bedrock (Nova)\\AWS Free Tier};
\draw[parrow, color=propred!80] (groq.east) -- (bedrock.west)
  node[above, font=\scriptsize, color=propred!80] {fail};

% Tier 3: Ollama (fork on Bedrock failure)
\node[tier, draw=awsgreen, fill=awslightgreen, right=2.0cm of bedrock]
  (ollama) {3. Ollama (Local)\\Inference API};
\draw[parrow, color=propred!80] (bedrock.east) -- (ollama.west)
  node[above, font=\scriptsize, color=propred!80] {fail};

% Tier 4: HuggingFace (fork on Ollama failure)
\node[tier, draw=propblue!70, fill=propblue!20, right=2.0cm of ollama]
  (hf) {4. HuggingFace\\Mistral-7B-Instruct};
\draw[parrow, color=propred!80] (ollama.east) -- (hf.west)
  node[above, font=\scriptsize, color=propred!80] {fail};

% Fallback (all fail)
\node[tier, draw=propred!80, fill=propred!20, right=2.0cm of hf]
  (fallback) {Fallback:\\Retrieval-Only};
\draw[parrow, color=propred!80] (hf.east) -- (fallback.west)
  node[above, font=\scriptsize, color=propred!80] {all fail};

% Success paths (arcs returning downward from each backend)
\node[rectangle, draw=awsblue, fill=awslightblue, rounded corners=4pt,
      minimum width=3.0cm, minimum height=0.7cm, below=1.2cm of bedrock,
      xshift=1.0cm, align=center, text width=2.8cm]
  (result) {Response Sent\\(with grounded citation)};

% Success arrows from each tier
\draw[parrow, color=awsgreen!80] ([xshift=0.5cm,yshift=-0.35cm]groq.south)
  arc (-90:0:0.5cm) -- ([xshift=-1.5cm]result.north)
  node[left, font=\scriptsize, color=awsgreen!80] {success};

\draw[parrow, color=awsgreen!80] (bedrock.south) -- (result.north);

\draw[parrow, color=awsgreen!80] ([xshift=-0.5cm,yshift=-0.35cm]ollama.south)
  arc (-90:-180:0.5cm) -- ([xshift=1.5cm]result.north)
  node[right, font=\scriptsize, color=awsgreen!80] {success};

\draw[parrow, color=awsgreen!80] ([xshift=-1.0cm]hf.south) -- ([xshift=2.5cm]result.north)
  node[right, font=\scriptsize, color=awsgreen!80] {success};

\draw[parrow, color=propred!70] 
  ([xshift=1.4cm]fallback.south) -- ([xshift=3.0cm]result.north)
  node[right, font=\scriptsize, color=propred!70] {degraded};

\end{tikzpicture}%
}
\caption{LLM backend fallback chain with automatic cascading.}
\label{fig:llm_fallback}
\end{figure}

\section{Deployment, Monitoring, and MLOps}
\label{sec:mlops}

Figure~\ref{fig:mlops_loop} presents the implemented MLOps lifecycle loop.

% ---- TikZ: MLOps lifecycle loop ----
\begin{figure}[H]
\centering
\resizebox{0.68\textwidth}{!}{%
\begin{tikzpicture}[node distance=1.4cm and 0.5cm]

% Main cycle nodes
\node[pbox=awsblue, minimum width=2.8cm] (ingest) {Data Ingestion\\(EventBridge)};
\node[pbox=awsgreen, right=2.0cm of ingest] (process) {Preprocess \&\\Embed};
\node[pbox=awspurple, right=2.0cm of process] (index) {Index\\(pgvector)};
\node[pbox=awsorange, below=1.8cm of index] (serve) {Serve\\(LLM + RAG)};
\node[pbox=propblue, below=1.8cm of ingest] (monitor) {Monitor\\(CloudWatch)};
\node[pbox=propgreen, below=1.8cm of process] (evaluate) {Evaluate\\(RAGAS + Stats)};

% Cycle arrows
\draw[parrow, thick] (ingest) -- (process);
\draw[parrow, thick] (process) -- (index);
\draw[parrow, thick] (index) -- (serve);
\draw[parrow, thick] (serve) -- (evaluate);
\draw[parrow, thick] (evaluate) -- (monitor);
\draw[parrow, thick] (monitor) -- (ingest);

% CI/CD overlay
\node[pbox=propgray, above=1.1cm of process, minimum width=5.0cm] (cicd)
  {CI/CD: GitHub Actions + SAM};
\draw[-{Stealth}, dashed, color=propgray, thick] (cicd.south west) ++(1.2,0) -- ++(0,-0.6) -| (process.north);
\draw[-{Stealth}, dashed, color=propgray, thick] (cicd.south east) ++(-1.2,0) -- ++(0,-0.6) -| (index.north);

% Budget guard
\node[pbox=propred, below=0.8cm of evaluate, minimum width=2.6cm] (budget)
  {Budget Guard\\(Cost Explorer)};
\draw[-{Stealth}, dashed, color=propred, thick] (serve.south) |- (budget.east);
\draw[-{Stealth}, dashed, color=propred, thick] (budget.west) -| (monitor.south);

\end{tikzpicture}%
}
\caption{MLOps lifecycle loop with continuous evaluation and monitoring.}
\label{fig:mlops_loop}
\end{figure}

\subsection{LLM Generation Configuration}

Generation is performed through the configured LLM backend with the following parameter bounds:

\begin{table}[H]
\caption{Implemented LLM Generation Parameter Configuration.}
\label{tab:llm_params}
\centering
\begin{tabular}{p{3.0cm}p{5.5cm}p{4.0cm}}
\toprule
\textbf{Parameter} & \textbf{Value} & \textbf{Rationale} \\
\midrule
Max Output Tokens & 300 & Cost control \\
Temperature & 0.2 & Reduce hallucination risk \\
Timeout & 15\,s (Groq) / 20\,s (Ollama, HuggingFace) / 15\,s read + 5\,s connect (Bedrock) & Tight budgets within 30\,s API Gateway limit \\
Retry Attempts & 0 (fallback chain provides redundancy) & Single attempt per backend; cascade on failure \\
Backoff Base & 1\,s (unused; retained for future configuration) & ---  \\
Context Window & 128k tokens (Groq, Bedrock) / 4096 tokens (Ollama) & Model-appropriate context \\
\bottomrule
\end{tabular}
\end{table}

If all LLM backends fail after exhausting the fallback chain, the system degrades gracefully by returning a retrieval-only summary with the fallback message ``The language model is temporarily unavailable. Below are the most relevant sources retrieved for your query.'' and logging the error to CloudWatch.

\subsection{Drift Detection and Remediation}
\label{sec:drift}

TechPulse implements a dual-criteria drift detection mechanism to identify degradation in retrieval quality over time:

\textbf{Probe query baseline:} At each drift detection cycle, 20 fixed canonical probe queries (stored in \texttt{probe\_queries.json}) are executed through the hybrid retrieval pipeline. The mean and standard deviation of cosine similarity scores across all probes are computed and persisted in the \texttt{drift\_baselines} database table.

\textbf{Detection criterion 1 --- Relative drop threshold:} The current mean similarity is compared against the most recent non-alert baseline. A drift alert is triggered if the relative drop exceeds 10\%.

\textbf{Detection criterion 2 --- Shewhart $3\sigma$ control chart:} The system retrieves the last 10 non-alert baselines via \texttt{\_get\_baseline\_history(n=10)} and computes the historical mean ($\mu$) and standard deviation ($\sigma$). The lower control limit is defined as:

\begin{equation}
  \text{LCL} = \mu_{\text{hist}} - 3\sigma_{\text{hist}}
  \label{eq:shewhart}
\end{equation}

A drift alert is triggered if the current mean similarity falls below LCL. This statistical process control approach detects systematic degradation that may not be captured by point-wise comparison alone. The Shewhart criterion requires a minimum of three stored non-alert baselines (\texttt{len(history) >= 3}); on fresh deployments with fewer than three baselines, only the relative drop threshold is active until sufficient history accumulates.

Both criteria are evaluated independently, and either triggering constitutes a drift alert. A human-readable explanation (e.g., ``relative drop 15.2\%; Shewhart $3\sigma$ breach (LCL=0.4321, value=0.3890)'') is included in the returned result, logged to CloudWatch, and routed to SNS notification.

\textbf{Remediation protocol:} The ordered response upon drift alert includes: (1) corpus audit to identify localized drift by source or topic cluster; (2) metadata filter re-calibration if off-topic documents enter the corpus; (3) re-indexing affected corpus segments with updated preprocessing; and (4) embedding model refresh (deferred to post-prototype phase).

\subsection{Health Monitoring}

A scheduled health-check Lambda runs every 5 minutes via EventBridge, probing four service dependencies:

\begin{enumerate}[leftmargin=0.5in]
  \item \textbf{Database connectivity} --- executes a lightweight test query against RDS PostgreSQL and reports document and chunk counts.
  \item \textbf{S3 reachability} --- verifies the data bucket is accessible.
  \item \textbf{SQS reachability} --- confirms the ingestion queue is operational.
  \item \textbf{LLM backend reachability} --- sends a minimal test prompt (``Respond with OK.'') to the configured LLM backend and reports which backend responded.
\end{enumerate}

Results are logged to CloudWatch. Three CloudWatch alarms are configured: RAG API errors ($>3$ in 5 minutes), Ingestion errors ($>2$ in 5 minutes), and DLQ depth ($\geq 1$ message), all routing to SNS for operational notification.

\subsection{CI/CD and Deployment Automation}

A four-stage CI/CD pipeline has been implemented using GitHub Actions:

\begin{enumerate}[leftmargin=0.5in]
  \item \textbf{Lint and Test} --- Python 3.13 setup, dependency installation, Ruff linting across \texttt{src/} and \texttt{tests/}, Pytest execution with coverage reporting.
  \item \textbf{Docker Build} --- Builds \texttt{techpulse-backend} and \texttt{techpulse-frontend} container images.
  \item \textbf{SAM Build} --- Builds and validates the SAM templates (\texttt{template-freetier.yaml} and \texttt{template.yaml}).
  \item \textbf{Deploy (main branch only)} --- Deploys to AWS via SAM with parameterized database credentials and VPC configuration from GitHub Secrets.
\end{enumerate}

All infrastructure is defined as code using AWS SAM, ensuring reproducible and consistent deployments. The SAM template defines: RDS PostgreSQL instance (\texttt{db.t3.micro}), S3 data and frontend buckets, SQS queue with DLQ, four Lambda functions (Ingestion, Preprocess, RAG API, Health Check), HTTP API Gateway, SNS alert topic, EventBridge schedules, and CloudWatch alarms.

\section{AWS Cloud Deployment Architecture}
\label{sec:cloud_arch}

Figure~\ref{fig:aws_deploy} presents the AWS cloud deployment topology, showing the physical service placement for the Free Tier configuration.

\begin{figure}[H]
\centering
\includegraphics[width=\textwidth]{aws_techpulse.png}
\caption{AWS cloud deployment architecture (Free Tier configuration).}
\label{fig:aws_deploy}
\end{figure}

\section{User Interface}
\label{sec:ui}

The user interface is implemented as a lightweight React single-page application. It provides natural language query submission, structured response rendering with citation display, and retrieval metadata transparency. All responses are constrained through the RAG pipeline; direct model invocation is not permitted. When retrieval evidence is insufficient or hallucination is detected, the system returns an explicit fail-safe abstention message. The sidebar maintains a recent-query history backed by browser \texttt{localStorage}; each history entry persists the full result object (answer, sources, metadata), enabling instant restoration of previous answers without re-issuing the API call.

\medskip
\noindent This chapter described the complete system architecture: the six-layer model, data engineering pipeline, document state lifecycle, preprocessing, storage schema, LLM fallback chain, drift detection, and cloud deployment topology. Chapter~5 presents the evaluation framework designed to measure retrieval and generation quality.
